% Lab 2: Dart and null safety
\documentclass[10pt,aspectratio=169,t]{beamer}
\usetheme[lab]{upbminimal}

\course{Application Development for Mobile Devices}
\decklabel{Lab 2}
\title{Dart and null safety}
\subtitle{Classes, null-safe types and collections}
\author{Dinu-Ștefan Rusu}
\institute{FILS, Universitatea Politehnica București}
\date{Week 4}

\begin{document}

\titleframe

\begin{frame}[eyebrow=Today]{What this lab is for}
  \vskip 2mm
  \begin{grid}[2]
    \cell[\faLock]{Null safety in code}{Reach for \code{?}, \code{?.}, \code{??} and \code{late} deliberately instead of by trial and error}
    \cell[\faLayerGroup]{Classes and inheritance}{Model a small domain with an abstract base class and two subclasses}
  \end{grid}
  \vskip 6mm
  \taskmeta{Time}{2 hours}{Submit}{Push to \code{admd-labs}, branch \code{lab-2}, before Lab 3}
\end{frame}

\begin{frame}[fragile,eyebrow=Setup]{Where your code goes}
  \begin{steps}
    \item Open the project from Lab 1 and create a branch:
\begin{shell}
git checkout -b lab-2
\end{shell}
    \item Create \code{lib/models/} directory. All classes live there.
    \item One class per file with \code{lower\_snake\_case} filenames: \code{user.dart}, \code{bank\_account.dart}, etc.
    \item Only \code{main.dart} holds widget code: just enough to display what your classes produce.
  \end{steps}
  \begin{hint}
    Dart convention: lowercase filename for lowercase class. The graders read your repository by these rules.
  \end{hint}
\end{frame}

\begin{frame}[eyebrow=Reference]{The null-safety toolkit}
  {\small
  \begin{columns}[T]
    \begin{column}{0.50\textwidth}
      \textbf{Types and operators}
      \begin{itemize}
        \item \code{String?}: may be null; plain \code{String} never can
        \item \code{?.}: read a member only if non-null
        \item \code{??}: use right-hand value if left is null
        \item \code{??=}: assign only if currently null
        \item \code{!}: assert non-null; throws if wrong
      \end{itemize}
    \end{column}
    \begin{column}{0.46\textwidth}
      \textbf{Declarations and promotion}
      \begin{itemize}
        \item \code{late}: assigned after construction
        \item \code{required}: caller must pass it
        \item \code{= 0.0}: default for optional parameter
        \item \code{is Book}: type test that promotes variable
        \item \code{x != null}: flow analysis promotes to non-null
      \end{itemize}
    \end{column}
  \end{columns}
  }
  \vskip 2mm
  \begin{takeaway}{Use \code{!} last, not first.}
    Every \code{!} you write is a promise the compiler can no longer check for you.
  \end{takeaway}
\end{frame}

\begin{frame}[fragile,eyebrow=Reference]{Null safety in one file}
\begin{dart}
class Room {
  final String code;            // always present
  final int? seats;             // may be unknown
  String? projector;            // may be absent
  late final DateTime opened;   // set later, read after

  Room({required this.code, this.seats, this.projector});

  String describeSeats() =>
      seats == null ? 'Unknown' : '$seats seats';
}
\end{dart}
  {\small
  \code{code} never null. \code{seats}, \code{projector} may be null, so compiler forces handling. \code{opened} has no value yet; reading it throws \code{LateInitializationError}.
  }
\end{frame}

\begin{frame}[eyebrow=Task I]{Users}
  {\small
  \begin{columns}[T]
    \begin{column}{0.55\textwidth}
      \begin{steps}
        \item Create \code{User} with \code{String name}, \code{int? age}, \code{String? email}.
        \item Make \code{name} required; leave \code{age}, \code{email} optional.
        \item Create three \code{User} instances with different null combinations.
        \item Display safely; show ``Not provided'' where value is null.
        \item Implement \code{getDisplayAge()}: return ``Age not specified'' or age as string.
      \end{steps}
    \end{column}
    \begin{column}{0.41\textwidth}
      \kv{Done when}{}
      \begin{checklist}
        \item Three users render, one with all optionals null
        \item Both branches of \code{getDisplayAge()} work
      \end{checklist}
    \end{column}
  \end{columns}
  }
\end{frame}

\begin{frame}[eyebrow=Task II]{Bank accounts}
  {\small
  \begin{columns}[T]
    \begin{column}{0.55\textwidth}
      \begin{steps}
        \item Create \code{BankAccount}: \code{late accountNumber}, required \code{holderName}, \code{balance = 0.0}.
        \item Add optional \code{bankBranch}.
        \item Implement \code{initializeAccount(String number)}.
        \item Add method to check if account number is set.
        \item Display accounts safely in ListView.
      \end{steps}
    \end{column}
    \begin{column}{0.41\textwidth}
      \kv{Done when}{}
      \begin{checklist}
        \item All accounts render with branch or ``N/A''
        \item Uninitialized account shows message
      \end{checklist}
    \end{column}
  \end{columns}
  }
\end{frame}

\begin{frame}[fragile,eyebrow=Scaffolding]{Just enough UI to see your objects}
\begin{dart}
ListView.builder(
  itemCount: accounts.length,
  itemBuilder: (context, i) {
    final a = accounts[i];
    return ListTile(
      title: Text(a.holderName),
      subtitle: Text(a.bankBranch ?? 'Not provided'),
    );
  },
)
\end{dart}
  \begin{takeaway}{Widgets are Lecture 3 and Lab 3.}
    Build visible rows only. \code{??} handles ``Not provided'' in widgets.
  \end{takeaway}
\end{frame}

\begin{frame}[eyebrow=Task III]{Library: model classes}
  {\small
  \begin{columns}[T]
    \begin{column}{0.55\textwidth}
      \begin{steps}
        \item Abstract \code{LibraryItem}: \code{title}, \code{author?}, \code{publishDate?}.
        \item \code{Book extends LibraryItem}: add \code{pageCount?}, \code{isbn?}.
        \item \code{Magazine extends LibraryItem}: add \code{issueNumber?}, \code{publisher?}.
        \item Pass shared fields via super constructor.
        \item Add \code{borrower?} (reuse Task I \code{User}) and \code{dueDate?} to \code{LibraryItem}.
      \end{steps}
    \end{column}
    \begin{column}{0.41\textwidth}
      \kv{Done when}{}
      \begin{checklist}
        \item Both subclasses compile without unset fields
        \item All optional fields use \code{?}
      \end{checklist}
    \end{column}
  \end{columns}
  }
\end{frame}

\begin{frame}[eyebrow=Task III]{Library: API and borrowing}
  {\small
  \begin{columns}[T]
    \begin{column}{0.55\textwidth}
      \begin{steps}
        \item \code{Library} holds \code{List<LibraryItem>}; add \code{addItem()} with validation.
        \item Implement \code{searchByTitle()}, \code{getItemsByAuthor()}, \code{getRecentItems()} with null handling.
        \item Add \code{borrowItem()} that checks if already borrowed.
        \item Add \code{returnItem()} that clears borrower and dueDate.
      \end{steps}
    \end{column}
    \begin{column}{0.41\textwidth}
      \kv{Done when}{}
      \begin{checklist}
        \item All methods handle null safely
        \item Borrowed items are protected
        \item Null queries return empty list
      \end{checklist}
    \end{column}
  \end{columns}
  }
\end{frame}

\begin{frame}[eyebrow=Troubleshooting]{Errors you will hit today}
  {\small
  \trouble{Non-nullable 'x' not assigned before use}{Give a value, use \code{?}, or mark \code{late}}
  \trouble{LateInitializationError on 'x'}{A \code{late} field read before it was set}
  \trouble{Can't access property on nullable receiver}{Use \code{?.}, \code{??}, or null check first}
  \trouble{Null check \code{!} used on null value}{Replace with \code{??} or real null check}
  \trouble{Type 'String?' can't assign to 'String'}{Promote with \code{if (x != null)} or \code{??}}
  }
\end{frame}

\closingframe{Push before you leave}{Branch lab-2 · one commit per task · running app is your evidence}

\end{document}
